\section{Conclusão}
\label{sec:conclusao}

A proibição da pulverização aérea no Ceará é uma política de saúde com cinco
anos de vigência plena, um vizinho que não a adotou e nenhuma avaliação. Este
projeto propõe avaliá-la com um desenho de fronteira na Chapada do Apodi,
organizado pela teoria da mudança da lei: cada elo da cadeia, do voo ao
nascimento, tem uma medida e um teste, e o efeito estimado só é atribuído à
hipótese da política depois de examinados os caminhos concorrentes.

O projeto também diz o que não conseguirá. O efeito esperado sobre a média
municipal é pequeno, e o desenho principal não o detecta. O que se pode
aprender é um intervalo para o efeito médio, um contraste mais fino pelo
calendário das gestações, o comportamento da perda fetal e, pela primeira vez,
a medida de quanto avião a lei removeu.

\subsection{Limitações}

\begin{itemize}
\item A exposição da gestante não é observada. Todos os efeitos são de
intenção de tratar no nível do município.
\item Há de dois a três municípios tratados no núcleo. A inferência depende de
métodos para poucos grupos tratados, que têm baixo poder.
\item O primeiro estágio depende de pedidos de acesso à informação, que podem
ser negados ou devolver registros incompletos.
\item A linha de base do método de aplicação é de 2006.
\item Um choque estadual simultâneo à lei só é fechado pelas especificações
secundárias.
\item Os desfechos econômicos não são cegos: a série estadual e os números de
dois municípios tratados já são públicos e foram vistos.
\end{itemize}

\subsection{Etapas}

\begin{table}[ht]
\centering\small
\begin{tabularx}{\textwidth}{@{}rL@{}}
\toprule
 & Etapa \\
\midrule
1 & Protocolar os pedidos de acesso à informação (Ministério da Agricultura, ADAGRI, SEMACE). \\
2 & Fixar as listas de municípios e assinar a pré-especificação. \\
3 & Montar o painel de nascimentos e óbitos fetais dos dois estados, 2005 a 2024. \\
4 & Calcular o efeito mínimo detectável do desenho principal e do desenho de calendário, com o pré-período. \\
5 & Medir a fração de domicílios próximos aos bananais. \\
6 & Estimar: desfecho primário, diferença tripla, calendário, desfechos secundários, mecanismos. \\
7 & Segundo ensaio: preços, custo-benefício e discussão do instrumento. \\
\bottomrule
\end{tabularx}
\end{table}

\subsection{Declarações}

\paragraph{Dados e código.} Todos os dados são públicos ou serão obtidos por
pedido de acesso à informação. O código está em repositório versionado.
Microdados não são redistribuídos.

\paragraph{Ética.} A pesquisa usa dados secundários, públicos e anonimizados, e
agregados por município. Não envolve contato com pessoas.

\paragraph{Conflito de interesses e financiamento.} O autor declara não haver
conflito de interesses. [Informar a fonte de financiamento, se houver.]

\paragraph{Uso de inteligência artificial.} Um assistente de inteligência
artificial foi usado na busca e leitura de fontes, na programação e na redação
de versões deste texto. As referências foram conferidas em registro
bibliográfico ou no próprio documento, e o autor revisou o conteúdo e responde
por ele.
